% Chapter 7: Testing
\chapter{Testing}

\section{Unit Testing}
Unit testing was performed on individual components to verify correct functionality.

\subsection{API Unit Tests}
\begin{table}[H]
    \centering
    \begin{tabular}{|l|l|l|l|c|}
        \hline
        \textbf{Test ID} & \textbf{Module} & \textbf{Input} & \textbf{Expected Output} & \textbf{Result} \\
        \hline
        UT-01 & API 1 - /predict/ & Cotton fabric image & Valid JSON with fabric details & PASS \\
        \hline
        UT-02 & API 1 - Input validation & Text file instead of image & HTTP 422 error & PASS \\
        \hline
        UT-03 & API 2 - Parameter calculation & Cotton (dirt=4) and Silk (dirt=2) & Temp=30°C, Spin=6min & PASS \\
        \hline
    \end{tabular}
    \caption{API Unit Tests}
    \label{tab:unit_api}
\end{table}

\subsection{Arduino Unit Tests}
\begin{table}[H]
    \centering
    \begin{tabular}{|l|l|l|l|c|}
        \hline
        \textbf{Test ID} & \textbf{Module} & \textbf{Input} & \textbf{Expected Output} & \textbf{Result} \\
        \hline
        UT-04 & calculatePumpTime() & 40 ml detergent & 1964 ms & PASS \\
        \hline
        UT-05 & runMotorCycle() & 1 min, Normal pattern & 5s CW, 2s pause, 5s CCW & PASS \\
        \hline
    \end{tabular}
    \caption{Arduino Unit Tests}
    \label{tab:unit_arduino}
\end{table}

\subsection{Linear Regression Validation}
Mean absolute error between predicted and actual detergent volumes: 0.72 ml (within acceptable tolerance).

\section{Integration Testing}
Integration testing verified that modules work together correctly.

\begin{table}[H]
    \centering
    \begin{tabular}{|l|l|l|c|}
        \hline
        \textbf{Test ID} & \textbf{Description} & \textbf{Steps} & \textbf{Result} \\
        \hline
        IT-01 & API 1 to API 2 pipeline & Send image → API 1 response → API 2 & PASS \\
        \hline
        IT-02 & API to Arduino communication & Send parameters via WiFi → Arduino receives & PASS \\
        \hline
        IT-03 & Full system end-to-end & Image capture → Analysis → Parameters → Hardware execution & PASS \\
        \hline
    \end{tabular}
    \caption{Integration Tests}
    \label{tab:integration}
\end{table}

% Figure for integration test log
\begin{figure}[H]
    \centering
    \includegraphics[width=0.9\textwidth]{img/c7diagram1}
    \caption{Successful End-to-End Integration Test Output}
    \label{fig:integration_log}
\end{figure}

\section{Performance Testing}

\subsection{API Response Times}
\begin{table}[H]
    \centering
    \begin{tabular}{|l|c|}
        \hline
        \textbf{Test Scenario} & \textbf{Average Response Time} \\
        \hline
        API 1 – Single image & 2.3 seconds \\
        \hline
        API 1 – Batch of 3 images & 5.8 seconds \\
        \hline
        API 2 – Single calculation & 0.15 seconds \\
        \hline
    \end{tabular}
    \caption{API Response Times}
    \label{tab:perf_api}
\end{table}

\subsection{Hardware Response Times}
\begin{table}[H]
    \centering
    \begin{tabular}{|l|c|}
        \hline
        \textbf{Operation} & \textbf{Actual Time} \\
        \hline
        Detergent pump activation & 12 ms \\
        \hline
        Relay switching & 8 ms \\
        \hline
        WiFi parameter reception & 1.2 seconds \\
        \hline
    \end{tabular}
    \caption{Hardware Response Times}
    \label{tab:perf_hw}
\end{table}

\subsection{Load Testing}
\begin{table}[H]
    \centering
    \begin{tabular}{|c|c|c|}
        \hline
        \textbf{Concurrent Users} & \textbf{API 1 Success Rate} & \textbf{API 2 Success Rate} \\
        \hline
        1  & 100\% & 100\% \\
        \hline
        5  & 100\% & 100\% \\
        \hline
        10 & 95\% (rate limited) & 100\% \\
        \hline
    \end{tabular}
    \caption{Load Test Results}
    \label{tab:load}
\end{table}
Rate limiting handled by tenacity retry mechanism.

\section{Usability Testing}
Five participants tested the mobile application.

\subsection{Task Completion}
\begin{table}[H]
    \centering
    \begin{tabular}{|l|c|c|}
        \hline
        \textbf{Task} & \textbf{Success Rate} & \textbf{Difficulty (1–5)} \\
        \hline
        Navigate to scan screen & 100\% & 1 (Very Easy) \\
        \hline
        Capture fabric image & 100\% & 2 (Easy) \\
        \hline
        Review analysis results & 100\% & 1 (Very Easy) \\
        \hline
        Start wash cycle & 100\% & 1 (Very Easy) \\
        \hline
        Use manual override & 80\%  & 3 (Moderate) \\
        \hline
    \end{tabular}
    \caption{Usability Task Completion}
    \label{tab:usability}
\end{table}

\subsection{Key Feedback}
\begin{itemize}
    \item “App is intuitive and easy to use” – 4/5 users
    \item “Results appear quickly” – 5/5 users
    \item “Manual override needs clearer labeling” – 3/5 users
\end{itemize}

\subsection{Improvements Made}
\begin{itemize}
    \item Added explanatory text in manual override screen
    \item Increased size of emergency stop button
    \item Added confirmation dialog for cycle start
\end{itemize}

\section{Bug Fixes and Improvements}

\subsection{Critical Bugs Fixed}
\begin{table}[H]
    \centering
    \begin{tabular}{|l|l|l|}
        \hline
        \textbf{Bug} & \textbf{Issue} & \textbf{Fix} \\
        \hline
        BUG-01 & API 1 returning incorrect JSON format & Added proper JSON parsing \\
        \hline
        BUG-02 & Arduino freezing during motor operation & Added watchdog timer \\
        \hline
        BUG-03 & Detergent pump running continuously & Added flyback diode \\
        \hline
        BUG-04 & API 2 returning negative water level & Added minimum validation (10L) \\
        \hline
        BUG-05 & iOS camera crash & Updated Expo Camera permissions \\
        \hline
    \end{tabular}
    \caption{Critical Bugs Fixed}
    \label{tab:bugs}
\end{table}

\subsection{Improvements Implemented}
\begin{table}[H]
    \centering
    \begin{tabular}{|l|c|c|}
        \hline
        \textbf{Improvement} & \textbf{Before} & \textbf{After} \\
        \hline
        Image size optimization & 5–10 MB & <1 MB \\
        \hline
        Connection pooling & New connection per request & Reused connections \\
        \hline
        Fabric analysis caching & None & 5‑minute cache \\
        \hline
        Arduino startup time & 10 seconds & 3 seconds \\
        \hline
    \end{tabular}
    \caption{Improvements Implemented}
    \label{tab:improvements}
\end{table}


\section{Testing Summary}
\begin{table}[H]
    \centering
    \begin{tabular}{|l|c|c|}
        \hline
        \textbf{Test Category} & \textbf{Tests} & \textbf{Pass Rate} \\
        \hline
        Unit Testing     & 15 & 100\% \\
        \hline
        Integration Testing & 8  & 100\% \\
        \hline
        Performance Testing & 12 & 92\% \\
        \hline
        Usability Testing   & 6 tasks & 97\% \\
        \hline
    \end{tabular}
    \caption{Testing Summary}
    \label{tab:summary}
\end{table}
All critical functionality validated. System ready for deployment.
